\subsection{Integral finiteness of the actual weight-zero algebra}
\label{sec:support-integral-finiteness}

The preceding finite-end theorem controls characteristic-$\ell$ curves.
This is enough to prove the required integral finiteness, including the
part of the deformation ring killed by powers of $\ell$.  We first
explain why the original local ordinary condition survives passage to
such a curve.  Throughout, $\mathcal O$ is a complete discrete valuation
ring of mixed characteristic $(0,\ell)$, $\ell\ne2$, with finite residue
field $k$ and uniformizer $\pi$.  Write $E=\operatorname{Frac}(\mathcal O)$.

\begin{lemma}[Specialization of the ordered ordinary flag]
\label{lem:support-integral-ordinary-flag}
Let $R$ be the original universal local deformation ring.  Its residual
representation is $\operatorname{diag}(\psi,1,\psi^{-1})$.  Let $I_t$ be
the intersection of the selected minimal primes of the radical of the
ordinary ideal saturated by the original element $\delta$.  Each selected
prime is contained in the kernel of the base point $R\to\mathcal O$.
Retain the original local character $\Psi$, the inertia group $I$, and
the multiplier-one subgroup $G^0$.  The quadratic character $\eta$
is unramified at $\ell$, so $I\subseteq G^0$.  The original character has $\Psi|_I=\varepsilon^e\tau$, with
$e\ge1$ and $\tau$ of finite order.  The original normalizing
element has the form $\delta=\ell\delta_0\delta_1$, where
$\delta_0,\delta_1$ are conjugation-invariant trace expressions.
Suppose $x\in G^0$ is in the local decomposition group and satisfies
$\overline\Psi(x)^2\ne1$.  The element $x$ need not belong to inertia.
These are the local conditions of the constructed datum\cite{integral-finiteness-source}.

Every local map $R/I_t\to A$, where $A$ is a complete local domain,
gives an ordinary representation over that same $A$: there is
$h\in1+\mathfrak m_A M_3(A)$ for which the multiplier-one matrices
are upper triangular and the ordered inertia characters are exactly
$\Psi,1,\Psi^{-1}$, mapped to $A$.  This includes characteristic
$\ell$ domains.
\end{lemma}

\begin{proof}
Here is the field-point argument, retaining the finite inertial
characters.  Put $H=I\cap\ker\eta\cap\ker\tau$, where $\tau$ is the
original finite local character.  This is normal in $G^0$.  The ordinary
ideal imposes, for all $s_i\in H$, the characteristic polynomials and
ordered products
\[
 \operatorname{charpol}(r(s))=
       \prod_{i=0}^2(X-\Psi(s)^{1-i}),\qquad
 \prod_{i=0}^2\bigl(r(s_i)-\Psi(s_i)^{1-i}I_3\bigr)=0.
                                                               \tag{IF1}
\]
The displayed order is part of the equations.  Its saturation contains
these equations.  No minimal prime of that saturation's radical
contains $\delta$: multiplication by $\delta$ is injective on the
quotient by the saturation, and also on its reduced quotient.
Since $\ell$ divides $\delta$, every selected component has
characteristic zero.

Fix a selected prime $P$ and set $D=R/P$, $L=\operatorname{Frac}(D)$.
Choose the original $t\in H$ for which $\Psi(t)$ has infinite order.
Its three prescribed eigenvalues are distinct.  Diagonalize $r(t)$ in
that order.  In (IF1), substitute successively $(s,t,t)$, $(t,s,t)$ and
$(t,t,s)$.  The first product fixes the first column, the last fixes
the last row, and the middle product fixes the middle diagonal entry.
Thus the matrices for $H$ are upper triangular with diagonal
$\Psi,1,\Psi^{-1}$.  The first common eigenspace and the kernel of the
last common eigenfunctional are unique.  Normality of $H$ and
conjugation invariance of its three characters make this ordered flag
$G^0$-stable.  Denote its three characters by $\kappa_i$.

We must still identify them on the full inertia group.  If $s\in I$,
some $s^n$ belongs to $H$.  Every eigenvalue of $r(s)$ is therefore
algebraic over $E$, and so is every coefficient of its characteristic
polynomial.  The base point gives $D\to\mathcal O$; localize
$D[1/\pi]$ at the kernel of the resulting map to $E$.  In this local
domain an element algebraic over $E$ and congruent to $a\in E$ is
equal to $a$: a nonzero algebraic element over a field is a unit,
whereas its difference from $a$ belongs to the maximal ideal.
Applying this observation coefficient by coefficient gives
\[
 \operatorname{charpol}(r(s))=
          (X-\Psi(s))(X-1)(X-\Psi(s)^{-1})\qquad(s\in I).
                                                               \tag{IF2}
\]
Apply the trace identity in (IF2) to $s,st,st^2$.  The resulting
Vandermonde matrix on $\Psi(t),1,\Psi(t)^{-1}$ is invertible.
Consequently $\kappa_i(s)=\Psi(s)^{1-i}$ for every $s\in I$.

The decomposition-group mark is sufficient to make this flag
integral; it need not be an inertia element.  Over the complete local
ring $R$, Hensel's lemma factors the characteristic polynomial of
$r(x)$ into distinct linear factors with roots $\alpha_i\in R$,
uniquely ordered by
$\overline\alpha_i=\overline\Psi(x)^{1-i}$.  Their differences are
units.  The spectral idempotents
\[
 e_i=\prod_{j\ne i}
               \frac{r(x)-\alpha_j I_3}{\alpha_i-\alpha_j}
                                                               \tag{IF3}
\]
belong to $M_3(R)$ and reduce to the three coordinate projectors.
The vectors $e_i\mathbf e_i$, rescaled to have $i$th coordinate one,
form a matrix $h\equiv I_3\pmod{\mathfrak m_R}$.  At the base point
these projectors are exactly the coordinate projectors and $h=I_3$.

On a selected component $D=R/P$, the generic ordered flag already
constructed is $x$-stable.  Since $r(x)$ has three distinct
eigenvalues, this flag is the flag of the columns of $h$ in one
of their six possible orders.  Each of these six flags is defined
over $D$.  Its $G^0$-stability and the prescribed scalar action of
inertia on its quotients are closed matrix equations.  They hold
over $\operatorname{Frac}(D)$ and hence over $D$.  Specialize along
the original map $D\to\mathcal O$ and then to $E$.  At that base
point the three inertia characters $\Psi,1,\Psi^{-1}$ are distinct,
because $\Psi(t)$ has infinite order.  Their ordered flag is unique,
so the permutation must be the identity.  This also explains
directly the role of selecting components through the base point.

It follows that the lower-triangular entries of $h^{-1}r(g)h$, and
the differences between its inertia diagonal and
$\Psi,1,\Psi^{-1}$, belong to every selected $P$.  They therefore
vanish modulo $I_t$.  Mapping this same matrix $h$ to $A$ proves
the assertion.  No residual inertia distinctness is used, and the
argument remains valid when $\ell=3$; indeed it also works for
non-domain local targets of $R/I_t$.
\end{proof}

\begin{lemma}[Global triangularity in characteristic $\ell$]
\label{lem:support-integral-global-triangularity}
Let $A$ be a local domain containing $k$, and let $r$ be an original
normalized rank-three representation over $A$ with residual diagonal
characters $\psi,1,\psi^{-1}$.  Suppose that on a finite-index subgroup
its multiplier-one characteristic polynomials equal those of the base
representation.  Then all its diagonal entries equal the corresponding
base entries, and it is globally upper triangular or globally lower
triangular, in the original coordinate ordering.
\end{lemma}

\begin{proof}
Write $G^0=\ker\eta$.  For $g\in G^0$, some positive power $g^n$
belongs to the given subgroup.  The eigenvalues of $r(g)^n$ are
algebraic over $k$, so the eigenvalues of $r(g)$ and its characteristic
polynomial coefficients are algebraic over $k$.  Reduction gives
their prescribed constants in $k$.  The local-domain observation
used above, now over $k$, proves
\[
 \operatorname{charpol}(r(g))=(X-\psi(g))(X-1)(X-\psi(g)^{-1})
                 \qquad(g\in G^0).                           \tag{IF4}
\]

Work temporarily over $F=\operatorname{Frac}(A)$.  Write $c$ and
$g_1$ for the original marks and $g^*=cg^{-1}c$.  Let $J$ be the
anti-diagonal matrix with anti-diagonal entries one.  Polarization and
the original normalization give
\[
 r(g^*)=Jr(g)^{\mathsf t}J,\qquad
 r(g_1)-r(g_1^*)=\operatorname{diag}(\zeta,0,-\zeta).
                                                               \tag{IF5}
\]
The residue of $\zeta$ is $\psi(g_1)-\psi(g_1)^{-1}\ne0$.
The three diagonal eigenvalues in (IF5) are therefore distinct.
Their coordinate projectors belong to the algebra generated by
$r(G^0)$.

We record an elementary replacement for a semisimplification argument.
If this representation were irreducible, those projectors and the
nonzero matrix entries would generate all of $M_3(F)$.  Indeed a
nonzero $(i,j)$ entry supplies the matrix unit $E_{ij}$ by multiplying
on either side by projectors.  A failure to reach every vertex from
every other vertex gives a proper invariant coordinate subspace;
otherwise products give all the matrix units.  Form the algebra
spanned by
\[
       \bigl(r(g),\psi(g),1,\psi(g)^{-1}\bigr)
                  \quad(g\in G^0)
          \quad\hbox{in }M_3(F)\times F^3.
\]
The image in $M_3(F)$ of the kernel of its projection to $F^3$ is a
two-sided ideal.  It cannot be zero, since that would make $M_3(F)$
commutative.  It is therefore all of $M_3(F)$.  But the trace equality
from (IF4), extended linearly, would then give
$\operatorname{tr}(E_{00})=0$, a contradiction.  This proof is valid
also when $\ell=3$.

There is consequently a proper invariant subspace.  Since it is
stable under the diagonal matrix in (IF5), it is a coordinate
subspace.  The symmetry of vanishing entries in (IF5) is
$(i,j)\leftrightarrow(2-j,2-i)$.  A subspace on coordinates
$\{0\}$ or $\{0,1\}$ makes all matrices upper triangular; one
on $\{2\}$ or $\{1,2\}$ makes them lower triangular.  The two
remaining cases, $\{1\}$ and $\{0,2\}$, force all four entries
joining the middle coordinate to the outer two to vanish.
In that case the middle diagonal entry is an eigenvalue, is algebraic
over $k$ by (IF4), and reduces to one, so is one.  The outer
two-dimensional representation has characteristic polynomial
$(X-\psi(g))(X-\psi(g)^{-1})$.  The same matrix-algebra argument,
using $\operatorname{diag}(\zeta,-\zeta)$, makes it reducible.
Its invariant coordinate line makes the full representation upper
or lower triangular.  All these zero entries descend from $F$ to $A$.

Each diagonal entry of a triangular matrix in (IF4) is algebraic over
$k$ and reduces to its prescribed base value, so equals that value.
Finally $r(c)$ has matrix part one.  Multiplication by $c$ transfers
all these assertions from $G^0$ to its other coset.
\end{proof}

\begin{proposition}[Finiteness from characteristic-$\ell$ curves]
\label{prop:support-integral-finiteness-criterion}
Give the variables of $P=\mathcal O[[z_1,\ldots,z_d]]$ integer
weights $w_i$.  Let $J\subset P$ be an ideal stable under taking
weight components, and put $S=P/J$.  Let $S_0$ be the image of the
weight-zero formal series.  Assume that every local homomorphism
$f:S\to A$ to a complete discrete valuation ring of characteristic
$\ell$ with finite residue field satisfies
\[
 f(z_i)=0\quad(w_i=0),\qquad
 \bigl[f(z_i)=0\text{ for all }w_i>0\bigr]\quad\hbox{or}\quad
 \bigl[f(z_i)=0\text{ for all }w_i<0\bigr].             \tag{IF6}
\]
Then $S_0$ is finite as an $\mathcal O$-module.  No reducedness or
$\mathcal O$-flatness hypothesis on $S$ is required.
\end{proposition}

\begin{proof}
The monoid $\{a\in\mathbb N^d:\sum_iw_i a_i=0\}$ is finitely
generated.  For completeness, its nonzero elements minimal under
coordinatewise comparison form a finite set by Dickson's lemma;
subtracting one such element and inducting on total degree proves
generation.  Let $b_1,\ldots,b_s$ be the corresponding nonconstant
monomials.  Every weight-zero formal series is a formal series in
these monomials.  This substitution converges: if $D$ is the largest
total degree of a chosen generator, a monomial of total degree $n$
uses at least $n/D$ generators.  Conversely every generator has
positive degree.  Thus the two adic filtrations are cofinal.
Consequently $S_0$ is a complete Noetherian local quotient of
$\mathcal O[[u_1,\ldots,u_s]]$, with maximal ideal generated by
$\pi,b_1,\ldots,b_s$.  When $s=0$ the conclusion is immediate.

We claim that every $b_j$ belongs to $\sqrt{\pi S}$.  Otherwise
choose a prime $Q\supseteq\pi S$ avoiding $b_j$.  In the complete
local domain $D'=S/Q$, the image $b$ of $b_j$ is nonzero and is a
nonunit.  Extend $b$ to a system of parameters.  Among the minimal
primes of the ideal generated by the other parameters, one avoids
$b$: otherwise that ideal would be primary to the maximal ideal,
contrary to the height theorem.  Quotient by such a prime.  Its
radical after adjoining $b$ is maximal, so the principal ideal
theorem shows that the resulting complete local domain has
dimension one.

A complete Noetherian local ring is excellent; normalization of this
one-dimensional domain is therefore finite\cite{integral-finiteness-stacks}.  It is a complete local
discrete valuation ring $A$ with finite residue field.  Equivalently,
one may take a local factor of the finite normalization; completeness
and the domain condition leave just one factor.  The composite
$S\to A$ is local, has characteristic $\ell$, and does not kill
$b_j$.  However a nonconstant monomial of total weight zero either
uses a weight-zero variable or uses variables of both signs.  Equation
(IF6) makes its image zero.  This contradiction proves the claim.

The weight-zero projection on $P$ is $\mathcal O$-linear and
preserves $J$.  It therefore gives an $S_0$-linear projection
$S\to S_0$ and proves the exact contraction
\[
                  (\pi S)\cap S_0=\pi S_0.             \tag{IF7}
\]
In particular, from $b_j^{n_j}\in\pi S$ we obtain
$b_j^{n_j}\in\pi S_0$, even when $S$ has $\pi$-torsion.
Thus $\mathfrak m_{S_0}\subseteq\sqrt{\pi S_0}$, and finite
generation gives $\mathfrak m_{S_0}^N\subseteq\pi S_0$ for some
$N$.  The quotient $S_0/\pi S_0$ has finite length over $k$.
The two adic topologies are now equivalent, so $S_0$ is also
$\pi$-adically complete.  Lifting a finite $k$-basis of its
quotient and successive $\pi$-adic approximation proves that these
lifts generate $S_0$ over $\mathcal O$.
\end{proof}

\begin{theorem}[The original integral finiteness assertion]
\label{thm:support-actual-integral-finiteness}
For each original admissible rank-three datum and its marks supplied
by the actual auxiliary-frame construction, let $U$ be any original
universal local ring and $S$ any associated normalized universal global
ring.  The $\mathcal O$-subalgebra of $S$ generated by weight-zero
products of its original representation entries is finite over
$\mathcal O$.  This is the literal assertion $\mathrm{FinA}$, with
all its quantifiers over $U$ and $S$\cite{integral-finiteness-source}.
\end{theorem}

\begin{proof}
Let $S$ be that global ring.  Its representation entries determine
maps out of $S$ by its universal property.  Their displacements from
the residual diagonal lift therefore generate the relative
cotangent space.  Choose finitely many of them, giving a surjection
$\mathcal O[[z_1,\ldots,z_d]]\to S$.  Conjugation by
$\operatorname{diag}(t,1,t^{-1})$ preserves the original normalizing,
away-from-$\ell$, and local component conditions.  For the last
condition use $t=1+T$.  The universal local representation gives an
automorphism of $R[[T]]$ fixing $T$ and reducing to the identity at
$T=0$.  The ordinary matrix equations are conjugated, while the traces
defining $\delta$ are unchanged.  Thus it preserves the saturated
ordinary ideal and its radical.  The minimal primes of the extended
radical are precisely $P[[T]]$ for its finitely many original minimal
primes $P$.  They are permuted by the automorphism; reduction at
$T=0$ shows that each is fixed.  The base representation is fixed as
well.  In particular the intersection defining $I_t$ is preserved.

Assign each chosen entry displacement its original matrix weight.
The kernel $J$ of the presentation is stable under weight components.
One can check this integrally using the universal unit $t=1+T$:
modulo any power of $\mathfrak m_S$, a relation has only finitely
many relevant weights, and its conjugates give a Laurent polynomial
in $1+T$ with those weight components as coefficients.  Substitution
$t\mapsto1+T$ is injective on Laurent polynomials over any ring.
Thus each component vanishes modulo every power of $\mathfrak m_S$,
and completeness proves the assertion.  The same argument gives
every unchosen entry its original weight.

Now let $f:S\to A$ be any curve tested in
Proposition~\ref{prop:support-integral-finiteness-criterion}.
Normalization and the original away-from-$\ell$ inertia equations
survive this local map.  Adic continuity follows by reduction modulo
powers of $\mathfrak m_A$.  Lemma~\ref{lem:support-integral-ordinary-flag}
supplies exactly the remaining ordered local condition, over $A$
itself.  This is therefore an original ordinary point.  The actual
auxiliary frame, finite-point theorem and
Theorem~\ref{thm:support-finite-end} give the base characteristic
polynomials on one finite-index subgroup chosen before $A$ and $f$.
Lemma~\ref{lem:support-integral-global-triangularity} now gives
(IF6).  The criterion proves that $S_0$ is finite over $\mathcal O$.

Every weight-zero product of original matrix entries belongs to
$S_0$.  The literal algebra in $\mathrm{FinA}$ is consequently an
$\mathcal O$-submodule of this finite module.  Since $\mathcal O$
is Noetherian, it too is finite.  This argument has used the original
global equation quotient throughout; it has not inverted $\pi$ or
discarded torsion from that quotient.
\end{proof}

Finiteness supplies an exponent for each fixed arithmetic datum.  The
proof does not bound that exponent, the length of the integral
augmentation conormal module, or their growth as the original fields,
characters and primes vary.  These quantitative requirements and
characteristic-zero two-sided global rigidity remain separate.
